\documentclass[10pt,a4paper]{article}
\usepackage[margin=22mm,headheight=14pt]{geometry}
\usepackage{fontspec}
\setmainfont{TeX Gyre Pagella}
\setsansfont{TeX Gyre Heros}
\setmonofont[Scale=0.86]{DejaVu Sans Mono}
\usepackage{amsmath,amssymb,mathrsfs,booktabs,tabularx,array,xcolor}
\usepackage{enumitem,fancyhdr,titlesec,microtype}
\usepackage[hyphens]{url}
\usepackage[colorlinks=true,linkcolor=blue!45!black,urlcolor=blue!45!black]{hyperref}
\hypersetup{pdftitle={Chen--Moriwaki: Source Issues and Formalization Repairs},
 pdfauthor={ArakelovGeometry workspace review},
 pdfsubject={Retrospective review of Chapter 3, sections 3.3--3.4}}
\definecolor{navy}{HTML}{173D58}
\definecolor{graytext}{HTML}{52616B}
\titleformat{\section}{\Large\sffamily\bfseries\color{navy}}{}{0pt}{}
\titleformat{\subsection}{\normalsize\sffamily\bfseries\color{navy}}{}{0pt}{}
\titlespacing*{\section}{0pt}{0pt}{12pt}
\titlespacing*{\subsection}{0pt}{12pt}{5pt}
\pagestyle{fancy}
\fancyhf{}
\fancyhead[L]{\small\sffamily Chen--Moriwaki \textbar\ Source issues and repairs}
\fancyhead[R]{\small\sffamily Chapter 3, \S\S3.3--3.4}
\fancyfoot[L]{\footnotesize\color{graytext}Manuscript v1.2 \textbullet\ Review dated 16 September 2026}
\fancyfoot[R]{\small\thepage}
\renewcommand{\headrulewidth}{0.3pt}
\setlength{\parindent}{0pt}
\setlength{\parskip}{5pt}
\setlist{nosep,leftmargin=18pt,topsep=4pt}
\setlength{\emergencystretch}{2em}
\newcommand{\code}[1]{{\small\nolinkurl{#1}}}
\newcommand{\loc}[2]{\textbf{Source location:} printed p.\,#1 / PDF p.\,#2.}
\newcommand{\hist}[1]{\par{\footnotesize\color{graytext}\textbf{Record:} #1}}
\newcommand{\Q}{\mathbb Q}
\newcommand{\R}{\mathbb R}
\newcommand{\A}{\mathcal A}
\newcommand{\E}{\mathscr E}
\newcommand{\PP}{\mathbb P}
\begin{document}
\thispagestyle{plain}
{\small\sffamily\color{graytext}ARAKELOVGEOMETRY \quad / \quad FORMALIZATION RETROSPECTIVE}
\vspace{10pt}
\begin{flushleft}
{\huge\sffamily\bfseries\color{navy}Chen--Moriwaki:\\[3pt]Source Issues and Formalization Repairs}\\[10pt]
{\large Chapter 3, sections 3.3--3.4}\\[5pt]
{16 September 2026}
\end{flushleft}

The Horizon formalization identified \textbf{a false degree argument in the proof
of Theorem 3.4.12} and \textbf{an unjustified cylinder/Borel identification in the
proof of Proposition 3.4.6}. It replaced the affected arguments while retaining
the intended extension and probability constructions. Several printed indices
and function domains also required correction. The findings below do not
refute the principal theorem conclusions.

This report records \textbf{two substantive proof issues}, \textbf{twelve groups
of notation or domain slips}, and \textbf{one supplementary clarification}.
It separately explains two problems in the original formal statements:
these must not be attributed to errors in the book.

\textbf{Source and pagination.} Huayi Chen and Atsushi Moriwaki,
\emph{Arakelov geometry over adelic curves}, version 1.2, 5 October 2019,
the user-supplied \code{adelic_curve.pdf}. The audited range is printed
pp.\,178--199, PDF pp.\,194--215; the start of \S3.5 on the last page is
outside scope. PDF numbers are \emph{one-based viewer page numbers}:
\[
\boxed{\text{PDF page} = \text{printed page} + 16.}
\]
These locators refer to this manuscript, not to a different edition.

\textbf{Reading guide, in priority order.}
\begin{center}
\small
\begin{tabularx}{\linewidth}{@{}lXrr@{}}
\toprule
ID & Finding & Printed p. & Report p.\\
\midrule
S1 & Invalid degree argument; finite-descent repair & 197--198 & \pageref{s1}\\
S2 & Cylinder versus Borel $\sigma$-algebras & 193 & \pageref{s2}\\
E1 & Finite cylinders versus general measurable functions & 198 & \pageref{e1}\\
T1--T4 & Local degree, integration subscripts, root-bound domains & 178, 186--188 & \pageref{t1}\\
T5--T8 & Finite extension, completions, evaluations, fibre cylinders & 189--192 & \pageref{t5}\\
T9--T12 & Measurability proof and tower formula indices & 194, 197--198 & \pageref{t9}\\
F1--F2 & Corrections to the formal integration statements & 194--198 & \pageref{f1}\\
\bottomrule
\end{tabularx}
\end{center}

\textbf{Evidence boundary.} This is a retrospective of Horizon's existing work,
checked against the rendered source pages, transcription comments, review records,
and the completed mathematical formalization. The repairs below are explained
in mathematical language, without references to formalization code.
No new formal proofs were written for this report.
The recorded final audit covers 998 formal declarations with no admitted
proofs. Verification details and the limits of the review are on
p.\,\pageref{evidence}.

\newpage
\section{S1. The degree argument in Theorem 3.4.12}\label{s1}
\loc{197--198}{213--214}
Proof of Theorem 3.4.12: on p.\,197, the paragraph introducing a prime
upper degree, the paragraph following the first displayed tower identity,
and the paragraph following the definition of $P(T)$.

\textbf{What goes wrong.} The printed proof reduces to prime
$[L_2:L_1]$ and, for $K_2\in\E_{L_2/K}$ not contained in $L_1$, puts
$K_1=K_2\cap L_1$. It asserts
\[
[K_2:K_1]\mid [L_2:L_1]
\quad\text{and hence}\quad [K_2:K_1]=[L_2:L_1].
\]
The divisibility is false for general subfields. Also, arbitrary finite
extensions do not automatically decompose into towers of prime-degree
extensions; the proposed reduction needs additional justification. Later,
primality of a polynomial's degree and $a\notin L_1$ are used to assert
that it remains minimal over $L_1$. Those facts alone do not ensure
irreducibility after extending the ground field.

\textbf{The counterexample recorded by Horizon.} Let $a=\sqrt[3]{2}$,
let $\zeta$ be a primitive cube root of unity, and take
\[
 K=\Q,\qquad K_2=\Q(a),\qquad L_1=\Q(\zeta a),\qquad
 L_2=\Q(a,\zeta).
\]
The two distinct cubic fields have intersection $\Q$, and $L_2$ is their
compositum, of degree six over $\Q$. Consequently
\[
 K_1=\Q,\qquad [K_2:K_1]=3,\qquad [L_2:L_1]=2.
\]
Thus even a \emph{prime upper degree} does not make the claimed divisibility
true. Moreover, $T^3-2$ is minimal for $a$ over $\Q$ but has the root
$\zeta a$ over $L_1$; the degree of $a$ over $L_1$ is two.
This exposes a proof defect, not a counterexample to tower compatibility.

\textbf{How Horizon repaired it.} The replacement descends a finite
extension to a sufficiently large finite base, rather than using the
intersection as that base. For algebraic $B/K$, finite $N/B$, and a
prescribed finite $F/K$ inside $N$, it constructs finite $A/K$ in $B$ and
finite $E/K$ in $N$ such that
\[
 A,F\subseteq E,\qquad B(E)=N,\qquad
 [E:A]=[N:B],\qquad [E:A]_s=[N:B]_s.
\]
Here $[\,\cdot:\cdot\,]_s$ denotes separable degree. The construction
descends a basis and its multiplication coefficients, together with the
separable part. Equality of separable degrees gives a restriction
bijection of embeddings into an algebraically closed extension.
Counting these embeddings yields the local weighted-sum identity;
dividing by the common separable degree gives compatibility of fibre
probabilities. The finite-cylinder tower identity follows, and uniqueness
of probabilities on the generated $\sigma$-algebra gives the general result.

\textbf{Why the replacement gives the required tower identity.}
\begin{enumerate}
\item Finite descent supplies a valid finite base change containing every
element needed for a prescribed cylinder function. In particular, the
multiplication map $E\otimes_A B\to N$ is an isomorphism.
\item The embedding bijection shows that restricting an upper place gives
exactly the lower probability distribution: the local separable degrees
are the multiplicities, and the global separable degree normalizes their sum.
\item Therefore the direct and iterated fibre integrals agree on finite
cylinders. Their probability measures have identical finite marginals.
\item Cylinder sets generate the specified measurable structure, so
uniqueness of finite measures identifies the direct and iterated
probabilities. Integrating gives the nonnegative tower formula.
\end{enumerate}
\hist{The false divisibility and this cubic-field counterexample were
explicitly recorded during Horizon's source transcription. The completed
formalization proves the finite-descent and probability-compatibility
arguments described above.}

\newpage
\section{S2. Cylinder and Borel \texorpdfstring{$\sigma$}{sigma}-algebras}\label{s2}
\loc{193}{209}
Proof of Proposition 3.4.6: the paragraph identifying
$\A_{L,\omega}$ with the Borel $\sigma$-algebra of the projective topology;
also the description of the extended measure near the end of the proof.

\textbf{The issue.} On the previous page, $\A_{L,\omega}$ is defined by
the finite-stage projections. Equivalently it is generated by finite
cylinders. The printed proof identifies this with the full Borel
$\sigma$-algebra of the inverse-limit topology, without a countability
argument. That identification is not valid for arbitrary uncountable
inverse systems.

\textbf{Why the general inference fails.} The source notes give the
standard test case $X=\{0,1\}^{I}$ for uncountable $I$, viewed as the
inverse limit of its finite-coordinate products. Every set in the
cylinder-generated $\sigma$-algebra depends on at most countably many
coordinates. A singleton is closed, hence Borel, but does not have that
property. Thus the two $\sigma$-algebras differ in this example.

\textbf{Qualification.} This example refutes the unrestricted
topological-to-measurable inference used in the proof. It is not being
presented as a constructed counterexample in the particular class of
adelic fibres. The recorded repair avoids the unsupported equality;
it does not prove that every adelic fibre has unequal $\sigma$-algebras.

\textbf{How Horizon repaired it.} It retains the compact inverse-limit
topology and uses it only where needed. Cylinder sets are clopen and
compact. If a disjoint sequence of cylinders has a cylinder as its union,
compactness gives a finite subcover, so only finitely many members are
nonempty. Finite additivity therefore gives countable additivity on the
cylinder algebra. Carath\'eodory extension then supplies a probability
measure on the \emph{specified cylinder $\sigma$-algebra}, with the
correct finite marginals. This establishes the measure construction
required by Proposition 3.4.6 and used downstream.

\textbf{The resulting mathematical object.} Writing $p_{F,\omega}$ for
restriction to a finite stage, the measurable structure remains exactly
\[
 \A_{L,\omega}
 =\sigma\bigl\{p_{F,\omega}^{-1}(B):F\in\E_{L/K},\;
 B\subseteq M_{F,\omega}\bigr\}.
\]
The probability of such a cylinder is the probability of $B$ at its
finite stage. The formalization also proves the homeomorphism with the
space of compatible finite-stage families, without identifying its full
Borel $\sigma$-algebra with the displayed one.
\hist{This distinction and the repaired Carath\'eodory construction are
explicit in Horizon's source notes and mathematical blueprint.}

\subsection{E1. The implicit cylinder reduction in the general tower case}\label{e1}
\loc{198}{214}
Proof of Theorem 3.4.12, paragraph beginning the remaining integral
identity in the general case, just before the display involving $L''=K_2L_1$.

The wording describes an arbitrary nonnegative measurable $f$ as if it
could be written $g\circ\pi_{L_2/K_2}$ for one finite $K_2/K$.
Such a factorization characterizes a cylinder function, not a general
measurable function. The intended reading is a reduction to cylinders,
as already explained by a monotone-class argument on p.\,197.
Horizon makes the passage explicit: first prove the cylinder identity,
then identify the two probabilities on the generating $\pi$-system of
finite-stage measurable cylinders, and finally integrate arbitrary
nonnegative measurable functions. This establishes the result without
assuming that every measurable function itself descends to a finite stage.
\hist{Supplementary clarification identified in this retrospective;
not counted as a third established mathematical error. The completed
Horizon proof already contains the required extension step.}

\newpage
\section{Printed notation and domain slips: T1--T4}\label{t1}
Each entry gives the literal problem, the intended correction, and the
existing formalization that uses the corrected objects. These are local
notation repairs; they do not require changing the result's conclusion.

\subsection{T1. The completion subscript in the local-degree sum}
\loc{178}{194} Equation (3.3), in the opening finite-separable construction.

The denominator field in the local degree is printed $K_w$, although
the fixed base parameter is $\omega$. The consistent identity is
\[
\sum_{x\in M_{K',\omega}}[K'_x:K_\omega]=[K':K].
\]
The letter $w$ here should be $\omega$. The formalization uses the
completion at the restriction of the extension absolute value. The
local-degree sum gives the global degree; dividing by that degree shows
that the fibre weights sum to one.
\hist{The printed $K_w$ was explicitly recorded during source transcription.}

\subsection{T2. Reversed extension in the logarithmic norm formula}
\loc{186}{202} Proof of Theorem 3.3.7(4), equation (3.10).

The left side uses $I_{K/K'}$, reversing the extension. It must be
\[
 I_{K'/K}(\ln f_b)=\frac{1}{[K':K]}\ln f_\lambda,
 \qquad \lambda=N_{K'/K}(b).
\]
The input is a function on $\Omega_{K'}$ and the output is a function on
$\Omega_K$. The formalized norm identity has precisely this direction.
Its proof takes logarithms of the local norm product formula and divides
by the extension degree, producing the normalized weighted fibre sum.
\hist{Flagged in the page transcription and explicitly corrected in the
mathematical blueprint.}

\subsection{T3. Root bounds stated on the base space instead of the extension}
\loc{187--188}{203--204} Proof of Theorem 3.3.7(4): the sentence after
(3.11), and the opening paragraph on p.\,188 introducing the archimedean bound.

On p.\,187, the bound for $\ln f_b$ by $g\circ\pi_{K'/K}$ is said to
hold on $\Omega_K\setminus\Omega_{K,\infty}$. Its domain should be
$\Omega_{K'}\setminus\Omega_{K',\infty}$. Likewise, the analogous
archimedean bound on p.\,188 belongs on $\Omega_{K',\infty}$, rather
than $\Omega_{K,\infty}$. The variable $b$ lies in $K'$.

The formalization states the bound for an absolute value $w$ on the extension field
whose restriction is the absolute value at a specified base point.
This keeps the base envelope and the extension evaluation distinct.
The integrable envelope is then pulled back to the extension space and
used to prove integrability of the logarithmic evaluation there.
\hist{The p.\,187 occurrence is flagged in the original transcription;
the parallel p.\,188 occurrence was additionally checked in this retrospective.}

\subsection{T4. The fields of the two normalized absolute values are reversed}
\loc{188}{204} Proof of Theorem 3.3.7(4), first prose sentence after
the opening displayed archimedean bound.

The normalized absolute values $|\cdot|_\omega^{1/\kappa}$ and
$|\cdot|_x^{1/\kappa}$ are described as being on $K'$ and $K$,
respectively. The fields must be \textbf{$K$ and $K'$}, in that order.
The root-bound proof treats the second absolute value as an extension
of the first. After normalization by the archimedean exponent, the
triangle inequality gives the required bound for a root of the minimal
polynomial over $K$.
\hist{Explicitly flagged in the source transcription; the corrected
roles are used in the formalized root-bound argument.}

\newpage
\section{Printed notation and domain slips: T5--T8}\label{t5}

\subsection{T5. A displaced prime in the parameter-space notation}
\loc{189}{205} Definition 3.4.1, sentence immediately after identifying
the two integration operators $I_{K''/K}$ and $I_{K'/K}$.

The parameter space appears once as $\Omega'_K$. The notation used
throughout the construction is $\Omega_{K'}$, where $K'$ is the
separable closure of $K$ in $K''$. Thus the sentence is about the
identification
\[
 \Omega_{K'}\simeq\Omega_K\times_{M_K}M_{K''}.
\]
This is a notational inconsistency, not a change in the underlying set.
The formalized finite extension parameter space is defined through the
separable part, with the purely inseparable extension handled canonically.
An absolute value on the separable part extends uniquely across the
purely inseparable extension, which is what makes the identification valid.
\hist{Visible in the original image and its faithful transcription;
identified as a separate notation slip in this retrospective.}

\subsection{T6. Two incorrect completion labels in a factorization argument}
\loc{189}{205} Proof of Lemma 3.4.2, the nonarchimedean paragraph
factoring the minimal polynomial $F$ and identifying a completion with a quotient.

With $K^{\mathrm{sc}}=K(\alpha)$, the polynomial should be factored in
\textbf{$K_v[T]$}, not in the printed $K_v^{\mathrm{sc}}[T]$.
For the factor selected by the extension $w$, the isomorphism is
\[
 K_w^{\mathrm{sc}}\simeq K_v[T]/(F_i),
\]
not the printed $K_v^{\mathrm{sc}}\simeq K_v[T]/(F_i)$.
Here $K_w^{\mathrm{sc}}$ is the completion of the separable subfield
at the restriction of $w$. The corrected proof factors over the base
completion and obtains the extension completion from the selected factor.

The formalization proves that the image of the completed separable
subfield is exactly the separable closure in the completed extension.
It then derives the required local/global separable-degree expression
for the finite-fibre weights.
The remaining extension of completions is purely inseparable, while
the completed separable part is separable over the base completion.
These two properties identify it as the maximal separable subextension.
\hist{Both printed labels are flagged in the transcription and corrected
explicitly in the mathematical blueprint.}

\subsection{T7. Evaluation functions in a finite tower have the wrong domain}
\loc{191}{207} Proof of Proposition 3.4.4, paragraph describing the
generators of $\A'_2$.

The text writes $(y\in\Omega_K)\mapsto|\alpha|_y$ for $\alpha\in K_2$.
The domain must be \textbf{$\Omega_{K_2}=\Omega_2$}. A base absolute
value does not by itself evaluate an arbitrary element of $K_2$.
The canonical measurable equivalence between the iterated and direct
finite extensions compares evaluations on the correct extension space.
Its measurability is proved by checking that the base projection and
the evaluations of elements of the extension field are preserved.
\hist{Flagged in the transcription and explicitly explained in the
blueprint's discussion of finite-tower evaluation domains.}

\subsection{T8. Fibre cylinder functions described as global functions}
\loc{192}{208} Last main paragraph, definition of $V_{L,\omega}$
immediately before (3.17).

The text calls these functions on $\Omega_L$, but writes them in the
form $f\circ(\pi_{L/K''}|_{M_{L,\omega}})$. The displayed composition
has domain \textbf{$M_{L,\omega}$}. Accordingly $V_{L,\omega}$ is the
space of fibre cylinder functions. The separate global space $V_L$
is introduced in Remark 3.4.7 on the next page.
The formalization distinguishes these two spaces. Closure under addition
and multiplication follows by passing two cylinder functions to a common
finite stage before combining them.
\hist{The domain mismatch is explicitly noted in the page transcription.}

\newpage
\section{Printed notation and domain slips: T9--T12}\label{t9}

\subsection{T9. The starting class in the measurability proof has the wrong domain}
\loc{194}{210} First sentence of the proof of Proposition 3.4.8.

The class $\mathcal H$ is initially described as bounded nonnegative
functions on $\Omega_K$, even though the functions are restricted to
$M_{L,\omega}$. They must be functions on \textbf{$\Omega_L$}.
The same paragraph subsequently uses $\Omega_L$ correctly.
The formalization takes the input on the algebraic extension parameter space and
proves that its nonnegative fibre integral is measurable on the base.
The proof starts with finite-stage functions and extends to the
generated measurable structure by a monotone-class argument and
monotone convergence.
\hist{Explicitly flagged in the page transcription.}

\subsection{T10. The point in the upper-fibre probability has the wrong letter}
\loc{197}{213} Proof of Theorem 3.4.12, final probability formula,
after the factors $Q_{i,j}$ and their associated points are introduced.

The preceding sentence names the points $y_{i,j}$, while the formula
uses $x_{i,j}$. In the intended factorization setup the mass formula is
\[
 \PP_{L_2,y}(\{y_{i,j}\})=\frac{\deg Q_{i,j}}{p}.
\]
The singleton braces here make the measure argument explicit; the
essential typographical correction is $x_{i,j}\rightsquigarrow y_{i,j}$.
The formalization assigns the singleton mass to a point of the actual
upper fibre, using its normalized local degree.
\hist{Explicitly flagged in the page transcription. This local notation
correction does not fix the degree problem S1.}

\subsection{T11. A missing sum and evaluation at the wrong point}
\loc{198}{214} First display, continuing the proof of Theorem 3.4.12.

The formula evaluates $g$ at $x$, although $x\in\Omega_{K_1}$ and
$g$ is a function on $\Omega_{K_2}$. It also drops the sum over $i$
in the middle expression. The intended calculation is
\begin{align*}
 (I_{L_2/L_1}f)(y)
 &=\sum_{i=1}^{r}\sum_{j=1}^{n_i}\frac{\deg Q_{i,j}}{p}\,g(x_i)\\
 &=\sum_{i=1}^{r}\frac{\deg P_i}{p}\,g(x_i)
 =(I_{K_2/K_1}g)(x).
\end{align*}
Both occurrences of $g(x)$ inside the weighted sums must be $g(x_i)$,
and the outer sum must remain in the middle expression.

This repairs the \emph{notation of the intended calculation}, conditional
on a valid compatible finite-stage setup. Because S1 invalidates the
printed choice of fields in general, these edits alone are insufficient.
After finite descent, Horizon proves that restriction $r$ pushes the
upper probability $P$ to the lower probability $Q$, and hence
$\int g\circ r\,dP=\int g\,dQ$. This gives the required integral comparison.
\hist{Both slips are explicitly flagged in the page transcription; see S1 for the
substantive replacement.}

\subsection{T12. The finite-stage indexing family has the wrong base}
\loc{198}{214} Proof of Theorem 3.4.12, general-case comparison of
$\A_2$ and $\A'_2$, sentence just before concluding their equality.

The last occurrence of $K_2\in\E_{L_2/L_1}$ in that paragraph should
be \textbf{$K_2\in\E_{L_2/K}$}. The paragraph began with finite
subextensions over $K$, which generate the direct extension's
$\sigma$-algebra. The formalization constructs the measurable equivalence using
those projective measurable structures and the correct base projections.
Passing from $K_2/K$ to $K_2L_1/L_1$ relates the two families of finite
projections and establishes measurability in both directions.
\hist{Explicitly flagged in the page transcription.}

\newpage
\section{F1--F2. Corrections to the original formal statements}\label{f1}
\textbf{These are formalization corrections, not additional errors in the
book.} The source's distinction between nonnegative integration and
signed integration supplies the correct mathematical contract.

\subsection{F1. Integrability must be tested with the nonnegative integral}
\loc{194--195}{210--211} Definition of signed fibre integration after
(3.18), and Proposition 3.4.9.

The source first defines the nonnegative integral with values in
$[0,+\infty]$. Signed fibre integration is defined where the integral
of the absolute value is finite; Proposition 3.4.9 then treats integrable
functions and the induced operator on $L^1$.

The original formal criterion tried to express integrability using a
real Bochner fibre integral. The formal integration convention gives that integral the value zero
when its integrand is not integrable. On an infinite fibre, this can
hide a divergent integral of $|f|$: a returned real zero is not evidence
of finiteness. Finite-fibre arguments cannot simply be reused here.

\textbf{The correction already made by Horizon.} For a projectively
measurable real function $f$, the criterion is
\[
 f\in L^1(\nu_L)
 \quad\Longleftrightarrow\quad
 \int_{\Omega_K}^{+}\!\left(\int_{M_{L,\omega}}^{+}|f(x)|\,
 d\PP_{L,\omega}(x)\right)d\nu_K(\omega)<\infty.
\]
Both integrals are nonnegative extended-real integrals, so divergence
is retained as $+\infty$ instead of disappearing into a default real value. Signed
disintegration is then proved under integrability, and the resulting
operator on $L^1$ is a continuous linear contraction.
\hist{This concern was already recorded before the September completion
run. Horizon proved the corrected criterion and the corresponding
signed disintegration and $L^1$ results.}

\subsection{F2. An unrestricted signed tower identity is too strong}
\loc{194--198}{210--214} Signed-integral convention on p.\,194,
Proposition 3.4.9 on p.\,195, and the nonnegative tower proof of
Theorem 3.4.12 on pp.\,197--198.

The pre-existing formal statement of algebraic tower compatibility demanded a
pointwise iterated-integral identity for \emph{every} real function,
without measurability or integrability. Totalizing a nonintegrable
Bochner integral at zero does not establish unrestricted Fubini.
This extra quantification was not justified by the source.

\textbf{The corrected theorem keeps the adelic equivalence and projection
compatibility}, and supplies the two appropriate integral statements,
under the canonical identification of the parameter spaces:
\begin{align*}
 I^+_{B/K}(I^+_{N/B}f)&=I^+_{N/K}f
 &&\text{pointwise, for measurable }f:\Omega_N\to[0,\infty],\\
 I_{B/K}(I_{N/B}f)&=I_{N/K}f
 &&\nu_K\text{-a.e., for globally integrable real }f.
\end{align*}
The signed version allows arbitrary almost-everywhere representatives.
The proof uses integrability on almost every fibre, the nonnegative
identities, and invariance under changes on null sets.
The corrected statement therefore expresses the book's integral
compatibility on its proper domain, while retaining the canonical
equivalence of adelic curves themselves.
\hist{Raised by Horizon's independent review on 16 September 2026;
the statement was corrected after comparison with the source and then
proved. The completed proof keeps the nonnegative and signed cases distinct.}

\newpage
\section{Evidence, provenance, and limits}\label{evidence}

\subsection{Source identity and method}
The source is the 466-page, user-supplied PDF of Huayi Chen and Atsushi
Moriwaki's \emph{Arakelov geometry over adelic curves}, version 1.2,
5 October 2019. The original and the registered reference copy are
byte-for-byte identical. All page references in this report use the
printed pagination and the corresponding one-based PDF pagination
of that document.

The original page images were inspected for every printed issue catalogued
here. The transcriptions deliberately preserve the printed slips;
the corrected interpretations reside in the mathematical blueprint and
the completed formalization.
No source PDF or source transcription was edited for this report.

\subsection{What happened during the formalization}
The historical review includes the source transcriptions, issue records,
mathematical blueprint, completion reports, and final verification
evidence. The relevant progression was:
\begin{center}
\small
\begin{tabularx}{\linewidth}{@{}lX@{}}
\toprule
Stage & Mathematical result or review finding\\
\midrule
Source review & Recorded notation slips, the degree counterexample, and the cylinder/Borel concern.\\
Domain review & Reconciled signed integration and the integrability criterion with the source.\\
Finite extensions & Proved measurability, logarithmic integrability, norm identities, and finite towers.\\
Algebraic measures & Constructed the cylinder probabilities and proved disintegration and properness.\\
Tower repair & Proved finite descent, compatible fibre probabilities, and algebraic tower integration.\\
Final review & Completed the remaining source constructions and checked all formal obligations.\\
\bottomrule
\end{tabularx}
\end{center}

The substantive corrections S1 and S2, and the formal-statement
corrections F1 and F2, are explicit in the historical records.
Most notation slips were preserved and flagged during page transcription.
The second occurrence in T3, the separate notation observation T5, and
the supplementary clarification E1 are identified here as retrospective
observations, rather than being assigned an unsupported discovery date.

\subsection{What the recorded verification establishes}
Horizon's final full build completed at 07:26:42 UTC on 16 September 2026,
without warnings or errors. The compiled inventory and a separate audit
checked 998 project declarations and found no admitted proofs or
nonstandard axioms. The proof system's ordinary foundations included
choice, propositional extensionality, and quotient soundness.
The source index covers all 26 entries for the targeted sections,
including the separate parts of Theorem 3.3.7.

These are existing Horizon checks, not new proof builds performed while
writing this report. Kernel checking certifies the formal statements under
their stated assumptions; identifying a printed typo and matching a
theorem to its source also requires mathematical and visual review.
The counterexamples and explanatory remarks quoted in substance here
are not claimed to be separately formalized errata theorems.

\textbf{Scope of the conclusion.} This is a catalogue of the issues found
in the recorded work and the present page review, not a claim that no
other error exists in the manuscript. No comparison with later editions
or author-issued errata was made. In particular, the review found no
separate defect in Proposition 3.4.13: its generated-$\sigma$-algebra
claim is established by proving that the measurable structure generated
by finite-stage projections equals the smallest $\sigma$-algebra making
the base projection and all absolute-value evaluations measurable.
This result is independent of the unsupported full-Borel identification
discussed in S2.

\end{document}
